% SPDX-License-Identifier: Apache-2.0
% covers: Switch
\datasheet{Switch}{A five-port switch of the network on chip, one per virtual channel.}

\dsfacts{\code{lib/parts/src/bus/noc/switch.rs}}{\code{txhdl\_parts::bus::noc::switch::Switch}}%
{\code{//docs:noc}}

\subsection*{Function}

A switch has a two-way link on each of its four sides and a fifth port, the exit,
where something that is not a node attaches. Its inputs are
\code{n\_in}, \code{s\_in}, \code{w\_in}, \code{e\_in} and
\code{x\_in}, and its outputs are \code{n\_out}, \code{s\_out},
\code{w\_out}, \code{e\_out} and \code{x\_out}. Each port passes one
\code{Pkt}, a beat of one AXI channel addressed to a node. The switch
routes in dimension order, X before Y. Its place in the lattice is two more
inputs, \code{col} and \code{row}, which whoever places it holds at
constants with \code{tie}, so every switch of a lattice is one type
(issue 635). A \code{Node} is two switches,
one per virtual channel.

\subsection*{Parameters}

\begin{center}\footnotesize
\begin{tabular}{@{}lp{0.7\textwidth}@{}}
\toprule
Parameter & Meaning \\
\midrule
\code{XB} & The width of a column coordinate, and of \code{col}, in
bits. A lattice is at
most \code{1 << XB} columns wide. \\
\code{YB} & The width of a row coordinate, and of \code{row}, in bits. A lattice is at
most \code{1 << YB} rows high. \\
\code{A} & The address width of the AXI link. \\
\code{D} & The data width of the AXI link. \\
\code{S} & The strobe width, which is \code{D / 8}. \\
\code{I} & The identifier width of the AXI link. \\
\bottomrule
\end{tabular}
\end{center}

The tables below are for \code{Switch<2, 2, 32, 32, 4, 2>}: two-bit
coordinates, on the link that \code{//soc} uses.

\dstables{Switch}

\subsection*{Behaviour}

\begin{itemize}
\item The switch holds one thing, \code{turn}, three bits: which
input the choice starts from. The port a packet leaves by holds
nothing, being a function of the packet's \code{dx} and \code{dy} and
of \code{col} and \code{row}; what needed a register is which input goes
when several could (issue 315).
\item The switch compares coordinates at \code{CW}, 8 bits. It sends a
packet east when \code{dx} is above \code{col}, and west when \code{dx}
is below it. In the right column, it sends the packet south when
\code{dy} is above \code{row}, and north when \code{dy} is below it. At its own node it sends the packet out of the exit.
\item Each cycle the switch works out, for every input with a packet,
the port that packet wants and whether that port is ready. An input
whose port is full does not hold up the others.
\item Of the inputs that can move, the switch takes one per cycle: the
first at or after \code{turn}, going round the five, and \code{turn}
becomes the one after whichever moved. So an input waits on at most
four others, where the fixed priority it replaced could starve one
outright: a busy link kept the exit from moving at all, since the exit
was last and the links were never all quiet.
\item The switch sends the packet it takes unchanged.
\end{itemize}

\subsection*{Verification}

\begin{itemize}
\item Two tests in \code{lib/parts/src/bus/noc.rs}, in
\code{//lib/parts:parts\_test}, drive one switch at 0,0 alone.
\code{a\_busy\_link}\allowbreak\code{\_and\_the\_exit\_take\_turns} has a link and the exit
offer every cycle for the east port, and asserts that each moves more
than fifty of two hundred and that the two differ by at most two.
\code{two\_inputs\_of\_a\_switch}\allowbreak\code{\_contend\_for\_their\_shared\_output}
measures two inputs leaving by one port, both offering every cycle and
one offering every third.
\item Eight more tests of the network in the same file run it inside
four nodes, and \code{a\_core\_draws\_a\_solid\_and\_a\_gpu\_fills\_it},
in \code{//soc:demo\_test}, runs a core, a GPU, a memory and a serial
port on a two by two lattice of nodes.
\item The switch's netlist is checked inside \code{Mesh}, two to a
node: \code{ex\_mesh} lowers a mesh of six nodes, and
\code{//docs:mesh\_sim\_mesh\_tb\_test} and \code{//docs:mesh\_vsim\_test}
run it under nvc and Verilator against the Rust run, at the mesh's
ports.
\end{itemize}

\subsection*{Used by}

\code{Node}, which holds two: \code{req} for requests and \code{rsp}
for responses.

\subsection*{Limits}

The switch moves one packet per cycle, not one per port.
\code{//docs:noc} measures the cost on the system of \code{//soc}: the
rasteriser draws its list in 36\,516 cycles on a direct link and in
78\,391 cycles across the lattice.
